\documentclass[openany]{ctexbook}
\usepackage{amsmath,amssymb}
\usepackage[hidelinks]{hyperref}
\newcommand{\note}[1]{\footnote{本批译注：#1}}
\title{代数几何基础第五卷预备笔记\\8.18：校订续译}
\author{据用户上传英译编辑稿译出}
\date{2026年11月15日}
\begin{document}
\maketitle
\chapter*{范围及校订状态}
本批覆盖底本第34页命题8.18的陈述、原稿证明及其后附记，止于以“$d\ge2$”结尾的省略号。下一处仍在第34页，为“It remains to explain”起首的未编号段落，并跨至第35页。此为独立增补，未并入原合订本。

本批保留原稿中维数、奇异性与断句等未决问题。8.8的缺失条件和相关证明缺口没有恢复；译校、编译检查不构成对这些数学断言的独立证明。
\chapter*{8.18\quad 命题（原稿陈述，含待核处）}
假设 $k$ 是代数闭域，$X$ 是维数 $d\ge2$ 的整概形（分别地，整且正规），并且在 $k$ 上固有。设 $T$ 是 $X$ 的有限闭子集，使得 $X-T$ 光滑，并设 $f|_{X-T}$ 非分歧。要使那些满足 $Y_\xi$ 几何不可约（分别地，几何整）且“维数为 $1$”的点 $\xi\in P^\vee$ 所成集合的补集余维至少为 $2$，原稿称下列条件为充要条件。\note{底本此处确实写“dimension 1”，而(a)、(b)和随后的证明均用 $d-1$。末尾附记又讨论是否应限于 $d=2$，故不能仅凭上下文断定是漏写 $d-$ 还是原拟限定曲面。正文明确引述这一未决读法，不把它作为已校定的普遍定理。这里设 $d=\dim X$，只是明示原稿随后使用的记号。}
\begin{enumerate}
\item[(a)] 对每个 $x\in T$，存在一个经过 $x$、维数为 $d-1$ 且不可约的超平面截面 $Y_\xi$，其中 $\xi\in P^\vee(k)$。\note{底本(a)印作 $P(k)$，但 $\xi$ 参数化超平面，命题首句及(b)均使用对偶空间；这里统一为 $P^\vee(k)$。下文作为参数空间出现的 $P$ 同样统一为 $P^\vee$。}
\item[(b)] $X'=X-T$ 对于 $f$ 是“直纹的”（原稿如此），或者存在 $X'$ 的一个超平面截面 $Y'_{\xi_0}$，其中 $\xi_0\in P^\vee(k)$，其维数为 $d-1$，具有〔“非”字待核〕奇异性，且不可约。\footnote{原脚注51：“字迹不清，询问 A.G.”。}\note{底本写“(non???) singular”。究竟要求奇异还是非奇异，现有文本不能确定；这里不择其一。“直纹的”是作者暂用术语，应连同8.17后关于该术语的自我修正一起理解。}
\end{enumerate}
\section*{原稿证明}
先假设 $X$ 几何正规。由8.14(a)，我们已经知道可以在 $P^\vee$ 中找到一个余维至少为 $2$ 的闭子集 $Z'$，使得对 $\xi\in P^\vee-Z'$，$Y_\xi$ 在 $k(\xi)$ 上可分且维数为 $d-1$。对于这样的 $\xi$，$Y_\xi$ 几何不可约与几何整是等价的，所以8.18所考察的两个问题〔底本夹注：“respé et non respé”，法文第50页，问号〕便是等价的。

另一方面，由5.6，使 $Y_\xi$ 几何整且维数为 $d-1$ 的点 $\xi\in P^\vee$ 所成的集合 $U$ 是开的。这里维数假设表达的是 $\phi_\xi$ 为 $X_{k(\xi)}$-正则。

我们将给出一个包含于 $U$ 的非空开子集 $P^\vee-Z$，其定义是直接的：取 $Z$ 为 $g(Y'^{\mathrm{sing}})$ 与那些由 $f(x)$、$x\in T$ 确定的 $P^\vee$ 中的超平面 $H_x$ 的并。\note{底本此处像集未标闭包；正文不增添横线。$X\in T$ 的大写 $X$ 订为点变量 $x\in T$，与同句 $f(x)$、$H_x$ 及(a)一致。}对于 $\xi\in P^\vee-Z$，$Y_\xi$ 光滑且维数为 $d-1$；又由6.1，它几何连通，因而几何整。

因此，我们要表明（解释、证明），$Z$ 的每一个余维为 $1$ 的不可约分支都与开集 $U$ 相交。然而，这些不可约分支就是各 $H_x$〔它们可能重复，但这并不要紧〕，以及当 $g(Y'^{\mathrm{sing}})$ 确有余维 $1$ 时的这个集合；后者对应于 $X'$ 对于 $f$“不是直纹的”这一情形。

（注意：我们使用了 $Y'^{\mathrm{sing}}$ 的不可约性。另一方面，要使后一个集合与开集 $U$ 相交，原稿称充要条件是 $g(\overline{Y'^{\mathrm{sing}}})$ 与 $U$ 相交，并在这里插入“它包含一个开稠密集”的说明。）\note{底本这第三处横线位于 $Y'^{\mathrm{sing}}$ 上方，与前两处记号不同；不能未经原稿核对将其改为像的闭包。底本这句括号没有正常闭合，“which contains an open and dense set”的指代不够清楚。正文保留其内容而不擅自补入一个新的拓扑断言。按通常可构造像的论证，可以考虑像在其闭包中包含开稠密子集，但这里不把这种解释认定为手稿缺文的恢复。}

这就证明了此情形下的8.18。如果不假设 $X$ 正规，则对 $X$ 的正规化应用前面的结果；论证是直接的，我在这里不写出细节。\note{这是作者的省略声明，不是本批完成了正规化步骤的全部假设核验；尤其原稿(b)的奇异性条件仍待核。}

注意：在〔respé，即“分别地”〕的情形，8.18包含在8.16中；更确切地说，要除去 $d=2$ 的情形。至于〔non-respé，即非“分别地”〕的情形，或许最好“不要求 $d=2$，而不只是 $d\ge2$”……\note{最后这句话按底本“not to require $d=2$ and not only $d\ge2$”保留，其否定与比较关系在当前转写中不清楚。不能据此确定应增加还是删除曲面限制，亦不能据此前文维数 $1$ 的歧义宣称已修复。}
\end{document}
